% Independent audit of Parts G--J.  Statements are intentionally restricted
% to what is proved below; no blanket completion claim is made.
\section{Spaces over a base: two decisive tests}

\subsection{The proposed interval realization does not work}
\paragraph{G.1 (Refuted).}
The proposed assignment
\[
 R(F)=\operatorname*{hocolim}_{I\in\mathcal I(\mathbb R)}F(I)\times I
       \longrightarrow\mathbb R,
\]
where $\mathcal I(\mathbb R)$ consists of bounded nonempty open intervals
and arrows are inclusions, need not have open-inverse-image homotopy
cosheaf $F$.  This failure already occurs for a represented cosheaf.
\begin{proof}
Let $F(U)=*$ if $0\in U$, and $F(U)=\varnothing$ otherwise. This is
the open-inverse-image cosheaf of the continuous map $\{0\}\to\mathbb R$.
For an explicit verification of codescent, if $0\notin U$, every term
of the covering Cech diagram is empty. If $0\in U$, the nonempty terms
are indexed by tuples of members of the cover which contain $0$;
their Cech nerve is contractible. Its homotopy colimit is $*$.

In the usual topological bar realization of the displayed homotopy
colimit, for any $x\in\mathbb R$ choose a bounded interval $I$
containing both $0$ and $x$. The summand $F(I)\times I$ of the
space of zero-simplices then supplies a point over $x$. Geometric
realization cannot remove this point. Taking $U=(1,2)$ shows that
$R(F)\times_{\mathbb R}U$ is nonempty, whereas $F(U)=\varnothing$.
These spaces cannot be weakly equivalent. The same obstruction holds
for a derived homotopy colimit in the slice, since its structure maps
include the indicated summands. This disproves the proposed formula;
it does not disprove essential surjectivity of the comparison functor.
\end{proof}

% Remaining material is deliberately separated from the geometric theorems.

\subsection{Open-test preservation exceeds shape submersions}
\paragraph{G.2 (Proved).}
The continuous maps whose strict pullback preserves all open-test weak
equivalences include maps which are not shape submersions in the sense
of \citet[Definition 3.23]{GHIJVolpe2025}. In particular, the projection from the
Cantor space $K=\{0,1\}^{\mathbb N}$ to a point is such a map.
\begin{proof}
Let $u:X\to X'$ be a weak homotopy equivalence. For every open
$U\subset K$, the restricted pullback is $u\times1_U:X\times U\to
X'\times U$. The singular-complex functor preserves products, and
products of Kan complexes preserve weak equivalences. Thus this is a
weak equivalence, including when $U$ is empty.

No nonempty open subspace $V$ of $K$ is essential. Indeed, if the
constant-sheaf functor $\Delta:\mathcal S\to\operatorname{Shv}(V)$
had a left adjoint, it would preserve products. Put $D_n=\{0,1\}$.
The coordinate functions $V\to D_n$ give a section of
$\prod_n\Delta D_n$. A section of $\Delta(\prod_n D_n)$ is locally
constant as a function to the discrete set $\prod_nD_n$. Its image
under the comparison to $\prod_n\Delta D_n$ would be the coordinate
family only if it were the inclusion $V\to K$ on underlying sets.
That inclusion is not locally constant: $K$ has no isolated points,
and the same holds for every nonempty open $V$. Hence $\Delta$ does
not preserve this countable product. This contradiction proves the
claim about essentiality. A shape submersion to a point requires a
basis of essential neighbourhoods, so $K\to*$ is not one.
\end{proof}

\paragraph{G.3 (Open).}
G.1 refutes the specific proposed realization, not the comparison
$\operatorname{Top}/\mathbb R[W_{\mathrm{open}}^{-1}]
\to\operatorname{CoShv}(\mathbb R;\mathcal S)$.
Neither a replacement realization nor full faithfulness is proved
here. Consequently the requested extension through intervals,
circles, graphs, manifolds, CW complexes and the listed non-locally
contractible spaces is also not settled. The dual raw-pullback
recognition theorem in G(3), an exact characterization of all
pullback-preserving maps in G(4), and the complete identification
of the raw-space theories with the stable duality in G(5) remain
open in this audit. G.2 refutes the proposed equality with Volpe's
shape submersions. G.5 proves the requested inclusion of all
local product projections in the preserving class.

\paragraph{G.4 (Proved).}
For a locally compact Hausdorff space $Y$ and a stable
infinity-category $\mathcal V$ admitting small limits and colimits,
Lurie's covariant Verdier duality identifies
$\operatorname{Shv}(Y;\mathcal V)$ with
$\operatorname{CoShv}(Y;\mathcal V)$, by
$A\mapsto(U\mapsto\Gamma_c(U,A))$
\citep[Theorem 5.5.5.1 and Proposition 5.5.5.10]{GHIJLurie2017}.
This known theorem does not by itself provide either raw-space
localization in Part G. Its displayed construction fails to be
an equivalence for pointed-space coefficients.
\begin{proof}[Proof of the unstable obstruction]
Take $Y=\mathbb R$ and the nonzero constant sheaf $A=S^0$ of
pointed discrete spaces. For a compact $K\subset U$, every
connected component of the open set $U$ meets $U\setminus K$:
a nonempty open interval cannot be a compact subset of
$\mathbb R$. A locally constant section which is the basepoint
outside $K$ is therefore the basepoint everywhere. All section
spaces here are discrete, so their pointed homotopy fibre
\[
 \Gamma_K(U,A)=\operatorname{fib}\bigl(A(U)\to A(U\setminus K)\bigr)
\]
is contractible. Its filtered colimit over compact $K$ is
contractible. Thus $\Gamma_c(-,A)$ is the zero cosheaf although
$A$ is not zero. The compact-support transform is not conservative
and cannot be an equivalence in this unstable setting.
\end{proof}

The attached parametrized-spaces theorem uses a different
homotopy theory: for a locally compact CW base it embeds the
usual parametrized homotopy theory as the locally homotopically
constant homotopy sheaves \citep{GHIJShulman2008}. It does not
assert that open-inverse-image tests present every cosheaf.
Its comparison employs spaces of sections and a fibrant
replacement; substituting raw open inverse images changes the
problem. This separates what is available in that source from
the unresolved realization in G.3.

\paragraph{G.5 (Proved).}
Every product projection $p:V\times U\to U$ preserves the
open-test class under strict pullback, with no local-contractibility
assumption on either factor. More generally, the preserving class
contains every map $p:Z\to Y$ admitting an open cover $\{Z_a\}$
such that $p(Z_a)=U_a$ is open in $Y$ and
$Z_a\cong V_a\times U_a$ over $U_a$. In particular it contains
all shape submersions in Volpe's sense, and G.2 makes that
containment strict.
\begin{proof}
We use open-cover homotopy codescent. If $T=\bigcup_a T_a$ is
an arbitrary open cover and $\Pi$ denotes weak homotopy type,
the augmented Cech diagram satisfies
\[
 \operatorname*{colim}_{[n]\in\Delta^{\mathrm{op}}}
 \coprod_{(a_0,\ldots,a_n)}
 \Pi(T_{a_0}\cap\cdots\cap T_{a_n})
 \simeq\Pi(T).
 \tag{G.5.1}
\]
No numerability or separation hypothesis is needed. Indeed the
singular simplices lying in some cover member form a simplicial
subset weakly equivalent to $\operatorname{Sing}(T)$:
subdivide finite sphere and disk representatives until their
simplices lie in cover members. For each such singular simplex,
the nonempty set of cover members containing its image has a
contractible Cech nerve. Applying this degreewise proves
(G.5.1). These are precisely the small-simplex and covering-sieve
statements of \citet[Lemma A.3.3 and Proposition A.3.2]{GHIJLurie2017}.

Let $u:X\to X'$ over $U$ be an open-test equivalence, and
write $X_A=X\times_U A$ for an open $A\subset U$.
For an arbitrary open $W\subset V\times U$, choose an open
rectangle cover
$W=\bigcup_a(V_a\times U_a)$.
For every finite tuple of indices, put
$V_J=\bigcap_{a\in J}V_a$ and $U_J=\bigcap_{a\in J}U_a$.
The corresponding intersections in the two pulled-back spaces
are exactly $V_J\times X_{U_J}$ and
$V_J\times X'_{U_J}$. Their comparison is a weak equivalence:
$X_{U_J}\to X'_{U_J}$ is one by hypothesis, and product with
any space preserves weak homotopy equivalences, as follows from
$\operatorname{Sing}(V_J\times X_{U_J})
=\operatorname{Sing}(V_J)\times\operatorname{Sing}(X_{U_J})$.
Thus the two Cech diagrams are termwise equivalent. Taking their
colimits in spaces and applying (G.5.1) proves the open test over
$W$. This proves the product-projection assertion.

Finally preservation is local on the source $Z$ of the base-change
map. Given any open $W\subset Z$, cover it by $W\cap Z_a$.
An open-test equivalence over $Y$ restricts to one over $U_a$,
since $U_a$ is open in $Y$.
Every finite intersection of these opens lies in one chart $Z_a$,
where the product argument just proved applies to arbitrary open
subsets. Consequently the pulled-back comparison is a weak
equivalence on every finite intersection. Formula (G.5.1) then
gives the equivalence over $W$. Since $W$ was arbitrary, this
establishes the assertion for the locally product map. The case
$K\to *$ used in G.2 is the product projection with $V=K$.
\end{proof}

\section{A second proof of the finite-poset realization theorem}

\paragraph{H.1 (Proved).}
Let $f:P\to Q$ be a proper functor of finite posets, meaning that
\[
 L(p,q)=\{r\in P:r\geq p,\ f(r)=q\}
\]
is weakly contractible whenever $f(p)\leq q$. The geometric direct
image along $g=|Nf|$ preserves least-vertex-constructible sheaves and
agrees with categorical right Kan extension, for space, spectrum,
and derived-module coefficients.

The proof below replaces BISITE D.11's initiality induction on the
poset of fibre cells by normalization of barycentric coordinates,
followed by geometric proper base change.

\begin{proof}
We first record the finite-poset coefficient calculation used below.
If $A$ is a finite poset and $D:A^{\mathrm{op}}\to\mathcal V$,
then
\[
 \Gamma(|NA|,\nu_A^*D)\simeq\lim_{A^{\mathrm{op}}}D.
 \tag{H.1.1}
\]
Here is a proof independent of BISITE D.7. Refine to the ordinary
open-simplex stratification of the order complex. Cellular sheaf
descent computes global sections as the limit over the face poset
$\mathcal C(A)$ of nonempty chains, ordered by inclusion. The
coefficient on a face $\kappa$ is $D(\min\kappa)$, and for
$\kappa\subseteq\kappa'$ the restriction is the map
$D(\min\kappa)\to D(\min\kappa')$. The cellular-descent
description follows by the open-star cover: a star has the
specified stalk at its central face, intersections give the
face incidences, and sheaf descent gives their limit. Equivalently,
the exit mapping space between two open simplices is contractible
when the first is a face of the second and empty otherwise;
straight-line contraction inside each simplex proves this
description. The functor
$\min:\mathcal C(A)\to A^{\mathrm{op}}$ is initial. Indeed
its comma category over $a$ is the face poset of $A_{\geq a}$,
whose nerve is the barycentric subdivision of the contractible
complex $|N(A_{\geq a})|$. Initiality proves (H.1.1), with its
canonical limit cone, for every complete coefficient category.

We next prove a topological lemma. Suppose $h:R\to[n]$ is a map
of finite posets satisfying the displayed outgoing-lift condition.
Choose $\lambda_i>0$ with $\sum_i\lambda_i=1$, and let
$Z_\lambda\subset |NR|$ be the inverse image of the point with
these barycentric coordinates. Set $R_0=h^{-1}(0)$.
Deleting the coordinates at positive levels and dividing those at
level zero by $\lambda_0$ defines a continuous map
\[
 r_\lambda:Z_\lambda\longrightarrow |NR_0|.
\]
If $R_0$ is empty, both spaces are empty. Otherwise this map has
weakly contractible fibres and is a weak homotopy equivalence.

We prove this assertion by induction on $n$. For $n=0$ it is the
identity. Suppose $n>0$ and take $z\in|NR_0|$ with positive support
the chain $\kappa_0$. Put $p=\max\kappa_0$ and
\[
 U_p=\{u\in R:u\geq p,\ h(u)>0\}.
\]
Give $U_p$ its level map to $\{1,\ldots,n\}\cong[n-1]$.
This map satisfies the same lifting condition, since a further
upper lift of $u\geq p$ remains in $U_p$. Its first fibre is
$L(p,1)$ and is weakly contractible. A direct coordinate calculation
identifies $r_\lambda^{-1}(z)$ with the geometric fibre of this
level map at weights $\lambda_i/(1-\lambda_0)$ for $i>0$:
the coordinates on $\kappa_0$ are fixed as $\lambda_0z$, and the
remaining coordinates range over chains of $U_p$. The induction
hypothesis makes this fibre weakly equivalent to $|NL(p,1)|$.
It is therefore weakly contractible.

All these spaces are finite polyhedra (the source is a finite
polytopal complex). The map $r_\lambda$ is proper. Proper base
change for sheaves of spaces \citep[Corollary 7.3.1.18]{GHIJLurie2009} identifies the stalk of
$r_{\lambda *}r_\lambda^*A$ at $z$, for any sheaf $A$, with the
sections of the constant sheaf $A_z$ on $r_\lambda^{-1}(z)$.
On a finite polyhedron $T$ those sections are
$\operatorname{Map}(\Pi T,A_z)$. They equal $A_z$ when $T$ is
weakly contractible. Thus the unit
$A\to r_{\lambda *}r_\lambda^*A$ is an equivalence on every
stalk, and hence an equivalence: sheaves of spaces on finite
polyhedra are hypercomplete, by finite homotopy dimension.
Applying this to each constant sheaf with value $K$ and then taking
global sections gives
\[
 \operatorname{Map}(\Pi|NR_0|,K)
   \simeq\operatorname{Map}(\Pi Z_\lambda,K)
 \quad(K\in\mathcal S).
\]
Yoneda proves that $r_\lambda$ is a weak equivalence and closes
the induction. This proof also shows the stronger unit statement
for every sheaf. Spectrum and derived-module versions follow by
testing on all shifts of the sphere or the rank-one free module;
these compact generators detect equivalences and commute with
stalk formation.

Now fix $x\in|NQ|$ with support $q_0<\cdots<q_n$. Apply the lemma
to the restricted poset $f^{-1}\{q_0,\ldots,q_n\}$. We obtain
$r:g^{-1}(x)\to|NP_{q_0}|$ and the unit equivalence just proved.
If $F:P^{\mathrm{op}}\to\mathcal V$ is a categorical coefficient,
its least-vertex sheaf restricted to $g^{-1}(x)$ is precisely the
inverse image along $r$ of the least-vertex sheaf attached to
$F|_{P_{q_0}}$. To check this without any fibre-cell initiality
argument, let $\nu_P:|NP|\to P$ be the minimum-vertex map to the
downset Alexandrov topology. The associated sheaf is $\nu_P^*F$.
The equality of minimum vertices gives the strict equality
\[
 \nu_P|_{g^{-1}(x)}
   = (P_{q_0}\hookrightarrow P)\,\nu_{P_{q_0}}\,r.
\]
Functoriality of inverse image gives the claimed identification.
Consequently
\[
 \Gamma(g^{-1}(x),F)
   \simeq \Gamma(|NP_{q_0}|,F|_{P_{q_0}})
   \simeq\lim_{P_{q_0}^{\mathrm{op}}}F.
\]

The categorical Kan value at $q_0$ is
$\lim_{(P_{\leq q_0})^{\mathrm{op}}}F$. The inclusion
$P_{q_0}^{\mathrm{op}}\to(P_{\leq q_0})^{\mathrm{op}}$ is initial:
its comma category at $p$ is $L(p,q_0)^{\mathrm{op}}$.
Quillen's Theorem A therefore identifies this limit with the one
just computed. This is the single-level categorical fibre formula;
the multi-level geometric comparison has already been proved by
the different argument above.

Finally $g$ is proper. The equality $\nu_Qg=f\nu_P$ supplies the
canonical mate $\nu_Q^*f_*F\to g_*\nu_P^*F$. At $x$, proper base
change identifies this mate with restriction of the Kan limit cone
to the fibre. Under the normalization just constructed it is the
initial-functor comparison in the preceding paragraph, and is an
equivalence. Stalkwise detection proves the claimed sheaf
equivalence. Its source is least-vertex constructible; this proves
preservation and identifies the actual adjunction mate, hence its
units and counits.
\end{proof}

\paragraph{H.2 (Refuted).}
The unmodified finite-poset assertion cannot be extended to all
finite categories with nonidentity idempotents while retaining one
least-vertex stratum for each object and all categorical presheaf
coefficients. The smallest possible endomorphism example already
fails the coefficient comparison.
\begin{proof}
Let $C$ have one object and two arrows $1,e$, with $e^2=e$.
The projection $p:C\to *$ is categorically proper: every category
over the terminal category is proper, since its outgoing lifting
categories have initial objects. Moreover $BC$ is contractible.
Indeed the zero $e$ gives a natural transformation from the
identity functor of $C$ to the constant functor at its object.

The least-vertex partition of $BC$ has one stratum, the whole
space. A constructible sheaf is therefore a local system, and
since $BC$ is contractible its global sections equal its stalk.
Take the categorical set coefficient $F(*)=\{0,1\}$ with
$F(e)$ the constant map to $0$. Its categorical right Kan
extension to the point is its set of fixed elements, namely
$\{0\}$. The one-stratum locally constant sheaf having the stated
stalk $\{0,1\}$ has two global sections. These cannot agree.
For a field $k$, the same example takes $F(*)=k^2$ and $F(e)$
the projection onto the first summand; the categorical limit is
$k$, whereas the geometric constant-sheaf sections are $k^2$.
There is no derived correction: $k[e]/(e^2-e)\cong k\times k$,
and the invariant-summand functor is exact.

The defect occurs before asking for a geometric properness
theorem: the category $C$ is not the exit category of its
one-stratum realization. Restricting coefficients to those which
invert $e$ removes this example, but is a substantive alteration
of the asserted extension. One object and two arrows is minimal
for the presence of a nonidentity endomorphism.
The program \texttt{code/GHIJ/small\_counterexamples.py}
checks the category laws, the outgoing lifting category, the
contracting zero transformation and the two-element action.
\end{proof}

\subsection{All outgoing lift categories for forgetting one mark}

We use Jansen's exit variance. Objects of $\mathcal C_{g,n}$ are
essential multicurves $\sigma$, and
\[
 \operatorname{Hom}(\sigma,\tau)
  =\{g:g\sigma\supseteq\tau\}/\Delta_\sigma.
\]
Here $\Delta_\sigma\cong\mathbb Z^\sigma$ is the multitwist
subgroup, acting on representatives on the right. Thus adding a
node points opposite to an exit arrow. This is precisely \citet[Construction 3.18]{GHIJJansen2023}; his Proposition 3.23 identifies the
opposite with the category of degenerations. His Corollary 6.6
and Theorem 6.7 give the exit/exodromy interpretation. The printed
inequality in the second assertion of Theorem 6.7 is a sign error;
the stability hypothesis throughout is $2g-2+n>0$.

Let $Q:\mathcal C_{g,n+1}\to\mathcal C_{g,n}$ forget the last mark
and stabilize, and put $F=Q^{\mathrm{op}}$ in entrance variance.
Write $q\sigma$ for the surviving multicurve, identifying repeated
images, and $\phi$ for the forgetful mapping-class homomorphism.

\paragraph{H.3 (Proved).}
Fix $\sigma$ upstairs and an exit arrow
$\beta:q\sigma\to\eta$ downstairs. Its outgoing lifting
category for $Q$ is either weakly contractible or equivalent to
$B\mathbb Z$. It is equivalent to $B\mathbb Z$ exactly when
forgetting the mark contracts a rational bridge of $\sigma$ and
$\beta$ smooths the node produced by that contraction. This
classification holds for every stable $(g,n)$.
\begin{proof}
Choose a representative $b$ of $\beta$ and a lift $\widetilde b$
through the forgetful mapping-class homomorphism. Conjugating
targets by $\widetilde b^{-1}$ reduces to
$\beta=[1]:\tau\to\eta$ with $\tau=q\sigma$ and
$\eta\subseteq\tau$. The forgetful homomorphism is surjective
by the Birman exact sequence (\citealp[Section 2]{GHIJKentLeiningerSchleimer2009}; the three-punctured sphere case is immediate).

An object of the outgoing lifting category is
$\alpha=[g]:\sigma\to\rho$ with $q\rho=\eta$ and
$\phi(g)\in\Delta_\tau$. Each curve of $\tau$ has at least one
preimage in $\sigma$, and a twist on such a preimage maps to its
twist. Therefore $\Delta_\sigma\to\Delta_\tau$ is surjective.
Multiplying $g$ on the right by a suitable element of
$\Delta_\sigma$ leaves $[g]$ unchanged and makes $\phi(g)=1$.
The vertical isomorphism $[g^{-1}]:\rho\to g^{-1}\rho$ then
identifies this lift with a normalized lift
$[1]:\sigma\to s$, where $s\subseteq\sigma$ and $qs=\eta$.

Set
\[
 H=\{d\in\Delta_\sigma:\phi(d)\in\Delta_\eta\}.
\]
The full category of normalized lifts has the exact presentation
\[
 \operatorname{Ob}=\{s\subseteq\sigma:qs=\eta\},\qquad
 \operatorname{Hom}(s,t)=
 \begin{cases}H/\Delta_s,&s\supseteq t,\\
                \varnothing,&\text{otherwise},\end{cases}
 \tag{H.3.1}
\]
with composition induced by addition in the free abelian group
$\Delta_\sigma$. Indeed the triangle condition is
$[h]\circ[1]=[1]$ as arrows with source $\sigma$, so its
representative must belong to $\Delta_\sigma$. Such twists fix
every submulticurve of $\sigma$, and the remaining verticality
condition is exactly $\phi(h)\in\Delta_\eta$. Quotienting a
morphism with source $s$ divides by $\Delta_s$. This proves
(H.3.1) on objects, morphisms and composition, not only on
components.

Only the component carrying the forgotten mark can become
unstable. If it does, it is a rational component with two
remaining special points. Contracting it has exactly one of two
effects: a bubble with a retained mark deletes one curve, or a
bridge identifies two curves. No second component becomes
unstable, since the numbers of special points of adjacent
components are unchanged by stabilization. Thus the twist map
is one of the following three coordinate maps:
\[
 \mathbb Z^\sigma\to\mathbb Z^\tau:
 \quad\text{a bijection of bases},\quad
 \text{deletion of one coordinate},\quad
 (a_1,a_2,\ldots)\mapsto(a_1+a_2,\ldots).
\]
In the first case there is one normalized object and its
automorphism group is trivial. In the deletion case let $s_{\max}$
contain all preimages of $\eta$ and the deleted curve. Then
$\Delta_{s_{\max}}=H$, so (H.3.1) makes it an initial object.
In the identification case, if the duplicate image belongs to
$\eta$, take both duplicate curves in $s_{\max}$; again
$\Delta_{s_{\max}}=H$ and the object is initial. If that image
does not belong to $\eta$, there is a unique normalized object
$s$, while
\[
 H/\Delta_s\cong
 \ker(\mathbb Z^2\xrightarrow{(a,b)\mapsto a+b}\mathbb Z)
 \cong\mathbb Z.
\]
Its generator is the difference of the two bridge twists. This
is exactly the stated $B\mathbb Z$ case. The reasoning is
independent of the genus and of automorphism groups of the
other components, since the triangle condition eliminated those
groups before the computation.
\end{proof}

\paragraph{H.4 (Proved).}
The stable-split enumeration gives $4,26,236$ boundary strata for
$\overline{\mathcal M}_{0,4}$, $\overline{\mathcal M}_{0,5}$ and
$\overline{\mathcal M}_{0,6}$ respectively. For forgetting
$5\to4$, all $47$ normalized stratum/smoothing cases comprise
$44$ contractible lifting categories and $3$ of type $B\mathbb Z$.
For forgetting $6\to5$, all $766$ such cases comprise $696$
contractible categories and $70$ of type $B\mathbb Z$.
\begin{proof}
The program \texttt{code/GHIJ/marked\_curve\_lifts.py} enumerates
each stable bipartition of the labelled set, quotients by swapping
its two sides, and enumerates every pairwise compatible family.
Compatibility is checked by emptiness of one of the four
intersections of sides. Such families are exactly the edge splits
of stable labelled genus-zero trees: one constructs the tree
successively by the compatible splits, and contracting all its
edges recovers the unique one-vertex tree. A family has at most
$n-3$ splits, so the enumeration is finite and complete.

Deleting the last label computes $q\sigma$ and detects deleted
and repeated images. Every subset $\eta\subseteq q\sigma$ is
enumerated. The program constructs all objects of (H.3.1) and
the rank of every nonempty morphism group; these groups have
no torsion because the twist constraints are disjoint coordinate
sums with coefficient one. It then either exhibits the initial
object or checks that the category has one object and automorphism
group $\mathbb Z$. Its reduced chain model is respectively zero
or the usual circle complex, giving integral homology through
degree three and the same Betti numbers over $\mathbb F_2$ and
$\mathbb Q$. The CSV and full group presentations are supplied.
These are finite presentations of the infinite lifting categories;
they are not falsely labelled finite quotient categories. Arbitrary
mapping-class representatives reduce to these cases by H.3.
\end{proof}

\paragraph{H.5 (Proved).}
For every stable $(g,n)$ the entrance forgetful functor
$F:T_{g,n+1}\to T_{g,n}$ is $W_\infty$-proper. It is smooth
only for $(g,n)=(0,3)$.
\begin{proof}
We give an algebraic proof of the contractibility needed for
properness, using graph-of-groups normal forms instead of lifted
geodesics. In exit variance, fix $\sigma$ and an incoming arrow
$\beta:\eta\to q\sigma$. Lifting a representative of $\beta$
and then normalizing the arrow upstairs identifies its incoming
lifting category with the poset
\[
 \mathcal P(\eta;\sigma)
   =\{\rho\supseteq\sigma:q\rho=\eta\},
 \qquad\text{ordered by reverse inclusion}.
 \tag{H.5.1}
\]
Here is the morphism check. A normalized object is
$[1]:\rho\to\sigma$. A triangle from this object to
$[1]:\rho'\to\sigma$ has representative in $\Delta_\rho$,
since both composites are the identity coset with source
$\rho$. Every such representative fixes $\rho$, and is the
same coset; hence there is exactly one triangle when
$\rho\supseteq\rho'$ and none otherwise. All normalization
isomorphisms are vertical because their downstairs representatives
belong to $\Delta_\eta$.

First take $\sigma=\varnothing$ and a nonempty $\eta$. Cut the
punctured surface $S$ along representatives of $\eta$. Van Kampen
gives a graph of groups with vertex groups $\pi_1(S_v)$ for the
complementary components and edge groups $\mathbb Z$ for the
cut curves. The edge injections are injective: each cut curve
is essential and each boundary loop in a negative-Euler-characteristic
surface is nontrivial. With $G=\pi_1(S)$, form the coset graph
whose vertices are $G/\pi_1(S_v)$ and whose edges are the
corresponding $G/\mathbb Z$ with the two incidence maps. The
normal-form theorem for this graph of groups says this graph is
a tree. Concretely, a maximal tree in the finite dual graph
reduces the presentation to amalgamated products and HNN
extensions; an edge path with no immediate reversal gives a
reduced word, and the amalgam/HNN normal form prevents a
nonempty such word from closing at its initial coset.

The nonperipheral part of (H.5.1) is the opposite face poset of
this coset tree. This description can be checked by point pushing:
a lift with one curve over each member of $\eta$ places the
extra point in a complementary component, and paths moving it
there are identified precisely modulo the corresponding vertex
group. A lift with a duplicated curve places it in the annulus
between the two preimages; paths are identified modulo that
curve's cyclic group. Deleting either of the duplicate curves
gives the two incidence maps. These assertions follow from
isotopy extension and the point-pushing exact sequence; they
are also the coset description underlying \citet[Section 7]{GHIJKentLeiningerSchleimer2009}. They do not require a choice of
hyperbolic metric or geodesic representatives.

A curve which becomes peripheral encloses the extra point and
one retained puncture. At most one can occur, and it cannot
coexist with a duplicated-image annulus. Deleting it attaches
one terminal edge to the barycentrically subdivided coset tree:
the other curves lift uniquely outside its twice-punctured disc.
Thus all such possibilities add leaves, without adding a cycle.
The nerve remains a tree and is contractible. If $\eta$ is
empty, the empty multicurve is terminal in (H.5.1), giving
contractibility directly.

For general $\sigma$, cut along its curves. On every component
not containing the forgotten point, missing curves of
$\eta\setminus q\sigma$ lift uniquely. In the component which
does contain that point, if it stays stable, the problem is
exactly the preceding one with retained boundary circles treated
as punctures. An added curve which becomes parallel to a fixed
boundary is one of the peripheral leaves just described. If the
component destabilizes, it is a pair of pants and admits no
essential refinement; keeping $\sigma$ already fixes the bubble
or bridge, so the remaining lift poset is a singleton. This
decomposes (H.5.1) into one of the contractible posets above and
singleton factors. All incoming lifts of $Q$ are therefore
contractible, proving properness of $F$.

For $(g,n)=(0,3)$ the target entrance category is terminal and
every functor to it is smooth. Every other stable moduli space
has a boundary node. Inserting a rational bridge carrying the
forgotten mark at that node gives the duplicate-curve case in
H.3; smoothing that node produces the lifting category
$B\mathbb Z$. Thus $F$ is not smooth in every other case.
\end{proof}

\paragraph{H.6 (Proved).}
For the universal stable curve
$\overline{\mathcal M}_{g,n+1}\to\overline{\mathcal M}_{g,n}$,
geometric direct image of constructible coefficients agrees with
right Kan extension on the entrance categories, for compactly
generated presentable coefficients. Point-fibre base change and
base change by inclusions of closed unions of boundary strata
agree with categorical proper base change. In positive genus
these are assertions about the moduli stack, not its coarse
space.
\begin{proof}
The forgetful map is the proper representable universal stable
curve. Over one stratum, normalize the nodal fibres. Its
components are proper smooth families of marked smooth curves
with disjoint node and marking sections. Ehresmann local
triviality, applied together with these sections, trivializes each
such family on a sufficiently small contractible analytic chart
of the stratum. Gluing the section pairs along the fixed dual
graph gives a trivialization of the stratified universal curve
over that chart. For positive genus one uses an orbifold chart
and retains the finite stabilizer action; the argument descends
along that chart. Proper base change now says that direct-image
stalks on the stratum are sections of one fixed stratified fibre
with locally constant transport. They are constructible.

The exit functor induced by the geometric map is $Q$. On
monodromy it forgets the point, giving $\phi$; on a degeneration
it erases the peripheral curve or identifies the two bridge
curves, giving $q$. These generate the morphisms in Jansen's
Construction 3.18. The relation given by a twist in the source
maps into the twist group of its stabilized image, so the map
on representatives descends and agrees with composition. The
natural exodromy equivalence therefore identifies inverse image
with restriction along $Q$, or along $F$ in entrance variance.
The now-proved preservation of constructibility allows both
right adjoints to restrict to these full constructible categories.
Their contractible uniqueness identifies the geometric direct
image with the categorical right Kan extension, including its
unit and counit.

For an actual point fibre, the stratified fibre is the stable
curve with its marked and nodal points distinguished. Its open
component strata are punctured surfaces of negative Euler
characteristic and hence are $K(\pi,1)$'s. Each puncture or
node branch has cyclic link group injecting into the appropriate
surface group. Thus its exit mapping spaces are the discrete
coset sets of these cyclic subgroups, with the expected surface
groups as open-stratum automorphisms. Point pushing identifies
these groups and cosets with the strict fibre of $Q$; the bubble
and bridge points have trivial automorphisms after their twists
are divided out. This gives the exit-category comparison for
the fibre itself, including its mapping sets, rather than merely
its total classifying space.

A closed union of boundary strata is a replete full cosieve in
entrance variance. Its inverse image is defined by the identical
object condition after applying $F$, and the morphisms are full;
the resulting square is a strict categorical pullback. This
statement concerns the actual boundary inclusion; in positive
genus a clutching normalization need not be an inclusion and is
not silently substituted for it. In these point and closed-boundary
squares the geometric and categorical comparisons are mates
of the same restriction identity under the identified adjunctions.
They agree, and are invertible by geometric proper base change
and H.5.

For compactly generated presentable coefficients, test sheaves
on the compact objects. Such tests preserve limits and filtered
colimits, hence commute with stalks, and are conservative.
Consequently the space-valued proper base-change and local
triviality arguments give the asserted coefficient version.
No inversion of stabilizer orders is required for this stack
statement.
\end{proof}

\paragraph{H.7 (Proved).}
The program \path{code/GHIJ/finite_poset_direct_image.py}
independently enumerates all
maps between one representative of every poset isomorphism class
with at most five objects, including the empty poset. The counts
of posets of sizes $0,1,2,3,4,5$ are $1,1,2,5,16,63$.
It examines $1,349,916$ monotone maps and recognizes $96,372$
proper maps. Properness is certified by explicit simplicial
collapses of every nonempty outgoing lifting order complex;
noncontractible cases are rejected by homology over
$\mathbb Q$ or $\mathbb F_2$. The program aborts rather than
classifying an acyclic but uncertified complex as contractible.
No such undecided case occurs in this range.

For every proper map and every nonempty target simplex, it
computes the categorical stalk and the actual geometric-fibre
cellular cohomology, for the constant coefficient and for each
principal-cosieve indicator. This gives $5,536,094$
stalk/coefficient cases, each computed over both $\mathbb Q$
and $\mathbb F_2$, with no discrepancy. A fibre cell is indexed
by a source chain meeting each level of the target simplex;
it is a product of simplices, of dimension
$|\kappa|-|\sigma|$. Its signed boundary deletes one vertex
while keeping every level nonempty, using the product-boundary
sign. The code verifies that consecutive boundary matrices
multiply to zero. Exact rational ranks use integer elimination
with gcd reduction, and binary ranks use bit arithmetic.
For a principal cosieve, extension from the corresponding
closed subcomplex reduces both computations to its intersection
with the relevant fibre or lower inverse image.

The supplied per-map CSV records all tests as successful; the
JSON gives the totals by source and target size. These numerical
comparisons check cohomology, not a replacement proof of the
canonical mate; the latter is established in H.1.

\paragraph{H.8 (Proved).}
By the independently reconstructed proof in D.10, the table agrees with every
entry of BISITE B.15. In the following table the smooth and
proper columns each apply to both $W_1$ and
$W_{D(\mathbb Q)}^{\mathrm{lc}}$; the last two columns distinguish
the intrinsic tests. Thus there are $48$ scalar comparisons and
no changes. Map labels refer to GALOIS D.3--D.5, as in BISITE.
\[
\begin{array}{l|cc|cc}
\text{map}&\text{smooth, both}&\text{proper, both}
 &\text{intrinsic }W_1&\text{intrinsic }\mathbb Q\mathrm{lc}\\\hline
\mathrm{D.3}:p&\text{yes}&\text{yes}&\text{yes}&\text{yes}\\
\mathrm{D.3}:q&\text{no}&\text{yes}&\text{no}&\text{no}\\
\mathrm{D.3}:qp&\text{no}&\text{yes}&\text{no}&\text{no}\\
\mathrm{D.4}:p&\text{yes}&\text{yes}&\text{no}&\text{no}\\
\mathrm{D.4}:q&\text{no}&\text{yes}&\text{yes}&\text{no}\\
\mathrm{D.4}:qp&\text{no}&\text{yes}&\text{no}&\text{no}\\
\mathrm{D.5}:H\to K&\text{yes}&\text{yes}&\text{no}&\text{no}\\
\mathrm{D.5}:q:R\to[1]&\text{no}&\text{yes}&\text{no}&\text{no}
\end{array}
\]
The proofs, reconstructed categories and program for these
entries are in Part D. No finite collection of pullback tests
is used as a substitute for the universal intrinsic condition.

\paragraph{Audit of the inherited proofs.}
BISITE D.11's fibre-cell initiality induction is complete once
its stated cellular-exodromy and finite-dimensional
hypercompleteness results are supplied; no contradictory case
was found. H.1 gives a different geometric argument. BISITE
F.15 agrees with Jansen after switching exit to entrance
variance; its correction of the printed stability inequality
is necessary. F.16's normalized lifting-poset and geodesic-tree
proof is complete using its cited point-pushing identification;
H.5 replaces the geodesic step by the amalgam/HNN normal-form
argument. F.17 correctly restricts its geometric comparison to
point fibres and the indicated actual closed boundary
inclusions; it does not prove arbitrary geometric pullback
compatibility. H.6 retains that scope in positive genus.

\section{The maximum diagonal-stable smooth class}

\paragraph{I.1 (Proved).}
Let $p:E\to B$ be a functor of ordinary small categories, and form
its strict diagonal $\Delta_p:E\to E\times_B E$. Then
\[
 p\text{ and }\Delta_p\text{ are }W_\infty\text{-smooth}
 \quad\Longleftrightarrow\quad
 p\text{ is a discrete fibration}.
\]
Consequently the discrete fibrations form the largest subclass of
$\operatorname{Liss}_{W_\infty}$ closed under diagonals. They are also
closed under identities, composition and strict pullback.
\begin{proof}
For smoothness use the incoming lifting category: given $x\in E$
and $a:b\to p(x)$, its objects are arrows $h:y\to x$ with
$p(y)=b$ and $p(h)=a$, and its arrows are the commuting triangles
over $1_b$. If $p$ is smooth, this category is weakly contractible,
and in particular nonempty.

Take any two such lifts $h:y\to x$ and $k:z\to x$. They give an
arrow $(h,k):(y,z)\to(x,x)$ in the strict pullback. A lift of this
arrow through $\Delta_p$ would be an arrow $l:w\to x$ satisfying
$(w,w)=(y,z)$ and $(l,l)=(h,k)$. Thus its incoming lifting category
is empty unless $y=z$ and $h=k$. Smoothness of $\Delta_p$ excludes
emptiness. Existence and uniqueness of each incoming lift is exactly
the discrete-fibration condition.

Conversely suppose $p$ is a discrete fibration. Its incoming lifting
categories are singletons. For an arrow $(h,k)$ into a diagonal
object, the discrete-fibration uniqueness gives $h=k$ and $y=z$;
the diagonal lifting category is again a singleton. Therefore both
maps are smooth. The usual unique-lift argument proves stability
under composition and strict pullback. Any diagonal-stable subclass
of the smooth maps satisfies the left side of the equivalence for
each of its maps and is therefore contained in this class.
\end{proof}


\paragraph{I.2 (Proved).}
In the presheaf Kan formalism with space coefficients, a discrete
fibration $p:E\to B$ has universally equivalent left and right Kan
extensions if and only if $p$ is an isomorphism. In particular the
stable classification by finite discrete bifibrations does not carry
over to unpointed spaces.
\begin{proof}
Pull back to any object $b\in B$. The fibre is a discrete set $I$.
Left and right Kan extension to a point are respectively the
coproduct and product functors $\mathcal S^I\to\mathcal S$.
On the family constantly equal to $*$ they yield the discrete set
$I$ and $*$, respectively. Hence equivalence forces $I$ to contain
exactly one element. In a discrete fibration with exactly one
object over each object of $B$, every arrow of $B$ has exactly one
lift with its specified target. This makes $p$ bijective on objects
and on all hom sets, and therefore an isomorphism. Conversely an
isomorphism has equal inverse, left-Kan and right-Kan functors,
universally under strict pullback.
\end{proof}


\paragraph{I.3 (Proved).}
Let $\mathcal V$ be a presentable closed symmetric monoidal
infinity-category with tensor product preserving small colimits in
each variable. The assignment
$\Phi(S)=\operatorname{Fun}(NS^{\mathrm{op}},\mathcal V)$ extends to
a lax symmetric monoidal correspondence functor with smooth right
legs, sending $X\xleftarrow aM\xrightarrow bY$ to $b_!a^*$.
This independently completes the coherence argument in BISITE F.10.
\begin{proof}
We use the precise universal property of spans in \citet[Corollary 2.6(2)]{GHIJCnossenLenzRamzi2025}: for an infinity-category $C$
and a pullback-stable wide subcategory $E$, restriction from the
span $(\infty,2)$-category with arbitrary left legs and right legs
in $E$ is an equivalence onto contravariant functors on $C$ whose
restrictions along $E$ have left adjoints and satisfy Beck--Chevalley.
The equivalence includes adjointable transformations and their
higher coherences. No diagonal or cancellation hypothesis occurs
in this statement.

Take $C=N\operatorname{Cat}$, so its pullbacks are the ordinary
strict categorical pullbacks. The incoming-lift characterization
of smoothness shows directly that smooth maps form a wide
subcategory closed under pullback and composition. Restriction of
presheaves is a coherent functor on $C^{\mathrm{op}}$, and all left
Kan adjoints exist by presentability. If $b:M\to Y$ is smooth,
the inclusion
$M_y^{\mathrm{op}}\to(y\backslash b)^{\mathrm{op}}$ is final:
its relevant comma categories are exactly the incoming lifting
categories. Hence
\[
 (b_!A)(y)\simeq\operatorname*{colim}_{M_y^{\mathrm{op}}}A.
\]
In every strict pullback square the fibre categories on the two
sides are identical. Applying this formula identifies the canonical
Beck--Chevalley mate as an equivalence. All hypotheses of the
quoted universal property now hold. It produces the correspondence
functor, and its value on a span is the stated composite.

For completeness the monoidal enhancement is also covered by
this universal property, and does not require choosing separate
associativity isomorphisms. For each $k$, external tensor is the
natural transformation
\[
 \prod_{i=1}^k\Phi(S_i)\longrightarrow
 \Phi\!\left(\prod_{i=1}^kS_i\right),\qquad
 (A_i)\longmapsto\boxtimes_i A_i.
\]
It is left-adjointable in every smooth variable. Indeed Fubini
for colimits and separate cocontinuity identify, with their
canonical units,
\[
 (\textstyle\prod_i b_i)_!(\boxtimes_i A_i)
       \simeq\boxtimes_i (b_i)_!A_i.
\]
Apply the transformation part of the span universal property to
$C^k$, with right class $E^k$. Products commute with strict
pullbacks, so both sides define functors on the corresponding
span category. The displayed transformation has a unique coherent
extension. The original associativity, symmetry and unit
coherences of external tensor lift through the fully faithful
restriction equivalence, including all higher relations. These
extensions are exactly a lax symmetric monoidal structure.

Right Kan adjoints to restrictions exist, $b_!$ has right adjoint
$b^*$, and pointwise tensor is closed by presentability and
separate cocontinuity. These are the remaining adjoints in the
2022 definition of a six-functor formalism. The projection formula
can also be checked directly by the fibre formula: restriction
of $b^*M$ to $M_y$ is constant at $M(y)$, and tensoring with this
object commutes with the displayed colimit.
\end{proof}


\paragraph{I.4 (Proved).}
The preceding proof does not follow from \citet[Proposition A.5.10]{GHIJMann2022} by taking $I$ to be all smooth maps and $P$ to be the
isomorphisms. That specialization violates the suitable-decomposition
cancellation requirement. Indeed $p:[1]\to*$ and $pg:*\to*$,
where $g$ selects $1\in[1]$, are smooth, but $g$ is not smooth:
the incoming arrow $0\to1$ has no lift. Thus $p,pg\in I$ does
not imply $g\in I$. The span universal property used in I.3 is
the needed replacement. Contractible uniqueness of individual
adjoints, without this extension statement or an explicit
unfurling construction, by itself does not state the necessary
coherence theorem. The program
\texttt{code/GHIJ/small\_counterexamples.py} enumerates the
incoming lifts for $p$, $pg$, and $g$, verifying this finite
cancellation counterexample.


\paragraph{I.5 (Proved).}
Opposition of the base categories transports I.3 to a correspondence
formalism with proper right legs and values
$\Psi(S)=\operatorname{Fun}(NS,\mathcal V)$. The maximum
diagonal-stable proper class is the discrete opfibrations. All
allowed maps of either Kan formalism satisfy $f^!=f^*$, so they
are cohomologically etale with trivial twist. These assertions
follow by applying $S\mapsto S^{\mathrm{op}}$ to I.1 and I.3;
incoming lifting categories become opposites of outgoing lifting
categories, and opposition preserves weak contractibility.


\paragraph{I.6 (Proved).}
For pointed-space coefficients, the analogous universally
ambidextrous discrete fibrations are exactly the inclusions of
unions of connected components, up to isomorphism over the base.
\begin{proof}
At a point fibre $I$, left and right Kan extension are wedge and
product. Applied to $S^0$ in every coordinate, their numbers of
components are $1+|I|$ and $2^{|I|}$. They differ when $|I|\geq2$,
including infinite cardinalities by Cantor's theorem. Thus each
fibre is empty or a singleton. Pull back along a base arrow
$0\to1$. A transition from a nonempty fibre over $0$ to an
empty one over $1$ gives the discrete fibration $\{0\}\to[1]$.
On a nonzero pointed coefficient its left Kan value at $1$ is
zero and its right Kan value is that coefficient. The reverse
emptiness pattern is excluded already by the existence of the
discrete-fibration reindexing function. Consequently the support
is a union of components, and all transitions there are bijections.
Conversely extension by zero on such a union is both Kan
extensions, universally after pullback.
\end{proof}


\paragraph{I.7 (Open).}
No proof or counterexample is supplied here for the reversible
accessible-localizer classification in Part I(2), including its
interval between integral homology equivalences and $W_0$.
For Part I(3), the unresolved assertion is an assignment
$\mathbb D\mapsto\mathbb D'$ on all derivators satisfying
Der~1--4, natural in exact morphisms, with
$W_{\mathbb D'}=W_{\mathbb D}^{\mathrm{lc}}$; we also do not
supply the alternative obstruction by two derivators equivalent
on finite posets but with different locally constant towers.
The weakest sufficient enhancement and an iteration theorem in
that generality remain unresolved. I.9 below proves a specified
localized-space realization and its finite iteration for the
enhanced categories of spectra and $D(\mathbb F_p)$, including
both second-stage calculations. It does not resolve the
bare-derivator assertion or the reversible-localizer classification.


\paragraph{I.8 (Proved).}
For spectra or $D(k)$, with $k$ a nonzero field, the universally
ambidextrous discrete fibrations are precisely the finite
discrete bifibrations: their fibres are finite sets, and all
reindexing functions along base arrows are bijections. Together
with I.2 and I.6 this gives the unstable and stable proper-map
classifications within the maximum diagonal-stable class.
Every map in that class is cohomologically smooth and etale,
with trivial dualizing twist, because $f_!$ is left adjoint to
$f^*$ and hence $f^!=f^*$.
\begin{proof}
At a point fibre $I$, ambidexterity equates coproduct and product.
An infinite product cannot preserve all coproducts: take the
family of diagrams with the coefficient unit in one coordinate
and zero in every other coordinate. The comparison is the
canonical map from the direct sum to the product of the
nonzero groups $\mathbb Z$ or $k$, and it is not surjective
for an infinite index set. Compactness of the sphere or of
$k$ detects this failure. Thus every fibre is finite.

Pull back along an arrow $0\to1$, and let
$r:I_1\to I_0$ be its reindexing map. If $i\in I_0$ has no
preimage, put a nonzero coefficient $M$ at $i$ and zero
elsewhere. Its left Kan value at $1$ is zero, while its right
Kan value contains $M$ as the sole nonzero factor. If $i$ has
$m\geq2$ preimages, use the same coefficient. Its right Kan
value is the homotopy limit of the star consisting of $m$
arrows $0\to M$, namely $(\Omega M)^{m-1}$; its left Kan
value is again zero. Taking $M$ to be the suspension of the
coefficient unit rules out this case. Hence every $r$ is
bijective. Conversely a finite discrete bifibration computes
both Kan extensions fibrewise, and the canonical finite
biproduct equivalence identifies them naturally under the
fibre-transport bijections and every strict base change.
\end{proof}

No invertibility of a prime is implicit in I.8. Its fibres have
no isotropy groups. If one returns to the larger smooth class,
the distinction matters: for $BC_p\to*$, trivial
$\mathbb F_p$ coefficients have homotopy orbits with nonzero
homology in every nonnegative degree and homotopy fixed points
with nonzero cohomology in every nonnegative degree. They are
not equivalent. The standard periodic resolution has
differentials alternating $1-g$ and $1+g+\cdots+g^{p-1}$;
both become zero on the trivial module. If $p$ is invertible
in a field, averaging splits invariants and the norm gives
ambidexterity for a finite $p$-group. These group projections
are not in the maximum diagonal-stable class: their diagonals
fail smoothness whenever the group is nontrivial. Thus this
diagnostic does not enlarge or qualify I.8.

\section{Strict two-categorical indexing: the ordinary-base case}

\paragraph{J.1 (Proved).}
Let $S$ be an ordinary small category, regarded as a locally
discrete strict $2$-category. On $2\operatorname{Cat}/S$ define
the presheaf fine class by requiring $s\backslash u$ to be a
Street weak equivalence for every $s\in S$. Its infinity-localization
is equivalent to $\operatorname{Fun}(NS^{\mathrm{op}},\mathcal S)$.
The dual over-slice class gives
$\operatorname{Fun}(NS,\mathcal S)$. These equivalences are
compatible with the Street-nerve comparisons to the contravariant
and covariant simplicial theories over $NS$.
\begin{proof}
Because the base has only identity $2$-cells, its lax under-slice
agrees with the strict $2$-comma object. Define
\[
 G(X\xrightarrow p S)(s)=s\backslash X,
 \qquad G(X):S^{\mathrm{op}}\longrightarrow2\operatorname{Cat},
\]
with strictly functorial reindexing by precomposition. For a
strict diagram $A:S^{\mathrm{op}}\to2\operatorname{Cat}$, define
its $2$-categorical Grothendieck construction $E(A)$ as follows.
Objects are $(s,x)$ with $x\in A(s)$; a $1$-cell
$(s,x)\to(t,y)$ is $(a:s\to t,h:x\to A(a)y)$; a $2$-cell
between such arrows exists only when their base arrows agree,
and is the corresponding $2$-cell between their $h$ components.
Strict composition in $A$ gives strict composition in $E(A)$.

For every $s$ there is a strict $2$-functor
\[
 R_s:s\backslash E(A)\longrightarrow A(s),
 \quad((t,y),a:s\to t)\longmapsto A(a)y.
\]
Its section takes $x$ to $((s,x),1_s)$. The arrows $(a,1)$
give a strict $2$-natural transformation from section followed
by $R_s$ to the identity. Thus $R_s$ is a Street weak
equivalence; a strict $2$-natural transformation supplies a
simplicial homotopy after the usual cylinder comparison. The
$R_s$ assemble naturally to $GE(A)\to A$.

There is also a natural strict $2$-functor over $S$,
\[
 j_X:X\longrightarrow EG(X),\qquad
 x\longmapsto(p(x),(x,1_{p(x)})).
\]
At $s$, the composite $R_s(s\backslash j_X)$ is the identity
on $s\backslash X$. For an object
$((t,(x,\alpha:t\to p(x))),\beta:s\to t)$ of
$s\backslash EG(X)$, the arrow with base $\alpha$ and
identity component gives a natural transformation to its image
under $(s\backslash j_X)R_s$. Consequently
$s\backslash j_X$ is a Street weak equivalence. Notice that
all displayed equalities are strict because the base is an
ordinary category; this is the precise place where that
hypothesis enters.

By definition $G$ preserves and detects the fine class. The
natural equivalence $GE\to1$ proves that $E$ sends pointwise
Street weak equivalences to fine equivalences. Together with
$j:1\to EG$, these are inverse equivalences after infinity-localization.
The projective model structure on strict $S^{\mathrm{op}}$-diagrams
in the combinatorial Thomason model category of strict
$2$-categories presents
$\operatorname{Fun}(NS^{\mathrm{op}},\mathcal S)$: apply the
projective diagram Quillen equivalence induced by
$c_2\operatorname{Sd}^2\dashv\operatorname{Ex}^2N_2$
\citep[Theorem 6.27 and Corollary 6.32]{GHIJAraMaltsiniotis2014}.
The standard strict-diagram comparison is
\citet[Proposition 4.2.4.4]{GHIJLurie2009}, applied on the
simplicial-set side of that Quillen equivalence.

To identify the comparison functor, observe the strict identity
\[
 N_2(s\backslash X)
    =N(S_{s/})\times_{NS}N_2X.
\]
It follows because the Street nerve preserves limits and the
base comma construction is the indicated strict pullback.
In the contravariant straightening theory over an ordinary
nerve, the value at $s$ is the weak homotopy type of this
pullback with the representable left fibration $N(S_{s/})$.
This is the co-Yoneda calculation for straightening: on a
simplex $\Delta^n\to NS$ the pullback has one contractible
component for each arrow $s\to\sigma(n)$, and the general
calculation follows by the simplicial resolution and homotopy
colimits. Thus the Street nerve has the same straightened
diagram as the equivalence already constructed. The over-slice
case follows by total $2$-categorical duality and invariance of
Street weak equivalences under that duality \citep[Theorem 6.6]{GHIJAraMaltsiniotis2022}.
\end{proof}

\paragraph{J.2 (Open).}
The proof of J.1 does not establish the requested recognition
for the walking noninvertible or invertible $2$-cell, the
one-object generating-$2$-cell examples, or a general strict
$2$-category. For those bases the strict comma identity used
above fails: lax triangles contain nonidentity $2$-cells, and
reindexing must retain their coherent action. No relative
strictification or straightening theorem with all the required
hypotheses has been proved here.

\paragraph{J.3 (Open).}
The universally pullback-preserving classes in Part J(2), the
complete boundary-piece tests, and the proposed lax-comma Der~4
base-change theorem in J(3) are not established here. The
absolute model structure and Theorem A do not imply those
claims. The verified statements available from the sources
are Theorem A for the specified slices
\citep[Theorem 5.2]{GHIJAraMaltsiniotis2018} and invariance
under every cell duality
\citep[Theorem 6.6]{GHIJAraMaltsiniotis2022}. Reversing only $2$-cells exchanges lax
and oplax coherence; reversing all positive-dimensional cells
exchanges the matched under/over fine theories. A full
four-variance base-change table requires the still-missing
recognition and lax-comma compatibility.

